\section*{Hoofdstuk \ref{ch:b1:species-interactions} --- Wisselwerkingen tussen soorten}

\begin{solution}{exo:b1:species-interactions:1}
\omterm{def:b1:species-interactions:types}{Mededinging} $-/-$ (twee eiken om licht); \omterm{def:b1:species-interactions:types}{predatie} $+/-$ (lynx, haas);
\omterm{def:b1:species-interactions:types}{parasitisme} $+/-$ (teek, hert); \omterm{def:b1:species-interactions:types}{mutualisme} $+/+$ (bij, bloem);
\omterm{def:b1:species-interactions:types}{commensalisme} $+/0$ (varen op bast); amensalisme $-/0$ (gras in de schaduw
van een eik).
\end{solution}

\begin{solution}{exo:b1:species-interactions:2}
Twee soorten die dezelfde beperkende hulpbron op dezelfde wijze gebruiken
kunnen in een stabiele omgeving niet onbeperkt samenleven. De vijf zangers
gebruiken hetzelfde voedsel maar in verschillende delen van de boom en op
verschillende tijden: hun nissen verschillen, dus geen twee van hen doen het
op ``dezelfde wijze''.
\end{solution}

\begin{solution}{exo:b1:species-interactions:3}
\omterm{def:b1:species-interactions:functional}{Hanteertijd}: de tijd om één prooi te vangen, te eten en te verteren, waarin
geen andere wordt genomen. \omterm{thm:b1:species-interactions:holling}{Aanvalstempo}: het per tijdseenheid doorzochte
oppervlak maal de vangkans. Grootste tempo $1/h = 2$ prooien per
uur.
\end{solution}

\begin{solution}{exo:b1:species-interactions:4}
Een soort waarvan de verwijdering de structuur van de \omterm{def:b1:ecosystem-organization:ecosystem}{levensgemeenschap}
verandert buiten verhouding tot haar talrijkheid, meestal doordat zij een
overheersende mededinger in toom houdt. Talrijkheid meet aanwezigheid, niet
uitwerking; een zeldzaam roofdier dat de aanstaande winnaar eet heeft een
uitwerking die zijn aantallen niet laten zien.
\end{solution}

\begin{solution}{exo:b1:species-interactions:5}
\omterm{prop:b1:species-interactions:lotka}{Isocline} 1: $N_1 + 0.5 N_2 = 500$ (snijpunten 500 op $N_1$, 1000 op
$N_2$); \omterm{prop:b1:species-interactions:lotka}{isocline} 2: $N_2 + 0.6 N_1 = 400$ (400 op $N_2$, 667 op $N_1$).
Zij kruisen elkaar met elke soort sterker door zichzelf begrensd ($\alpha = 0.5 <
K_1/K_2 = 1.25$; $\beta = 0.6 < K_2/K_1 = 0.8$): stabiel samenleven
bij $N_1 = 429$, $N_2 = 143$. Met $\alpha = 2$: \omterm{prop:b1:species-interactions:lotka}{isocline} 1 snijdt bij
500 en 250, beide binnen die van \omterm{prop:b1:species-interactions:lotka}{isocline} 2 (667 en 400): soort 2
wint altijd.
\end{solution}

\begin{solution}{exo:b1:species-interactions:6}
$f = aN/(1 + ahN)$: bij 10, $2/(1 + 0.5) = 1.33$ per uur; bij 100,
$20/(1 + 5) = 3.33$ per uur (maximaal $1/h = 4$). Halve verzadiging bij
$N = 1/(ah) = 20$ per vierkante meter.
\end{solution}

\begin{solution}{exo:b1:species-interactions:7}
Bij een respons van type III stijgt het deel van de prooi dat wordt gegeten
met de dichtheid bij een lage dichtheid: zeldzame prooi wordt naar
verhouding minder genomen en kan zich herstellen. Bij een respons van type
II is het gegeten deel het grootst bij een lage dichtheid (het roofdier is
daar nooit verzadigd), zodat een kleine prooipopulatie verder omlaag wordt
geduwd.
\end{solution}

\begin{solution}{exo:b1:species-interactions:8}
\emph{P.~caudatum} gebruikte dezelfde bacteriën in dezelfde waterkolom
als \emph{P.~aurelia}, en de snellere, zuinigere \emph{P.~aurelia}
drukte het voedsel onder wat zij nodig had. \emph{P.~bursaria} voedde zich
op de bodem met andere hulpbronnen: de \omterm{def:b1:ecosystem-organization:niche}{nis}, niet de groeisnelheid, gaf de
doorslag --- \emph{P.~bursaria} groeit niet sneller dan \emph{P.~caudatum}.
\end{solution}

\begin{solution}{exo:b1:species-interactions:9}
Roofdieren komen het patroon vaker zonder schade tegen, leren het minder
goed en vallen het vaker aan: de bescherming daalt voor zowel de nabootser
als het voorbeeld. De fitness van de nabootser daalt dus naarmate hij
talrijker wordt, een negatieve frequentieafhankelijkheid die hem zeldzaam
houdt ten opzichte van het voorbeeld.
\end{solution}

\begin{solution}{exo:b1:species-interactions:10}
Zonder het roofdier ligt de \omterm{prop:b1:species-interactions:lotka}{isocline} van de mossel buiten die van elke
mededinger: zij wint de rots. De \omterm{def:b1:species-interactions:types}{predatie} verlaagt de effectieve $K$ van de
mossel (haar \omterm{prop:b1:species-interactions:lotka}{isocline} wordt naar binnen getrokken) tot zij de \omterm{prop:b1:species-interactions:lotka}{isoclinen} van
de anderen van binnenuit kruist: de mededingers worden nu sterker door
zichzelf dan door de mossel begrensd, en leven met haar samen. Het roofdier
werkt op de sterkste mededinger in en maakt plaats voor de rest.
\end{solution}

\begin{solution}{exo:b1:species-interactions:11}
Wanneer fosfaat de groei begrenst is \qty{80}{\%} van het fosfaat veel meer
waard dan \qty{15}{\%} van de suiker: \omterm{def:b1:species-interactions:types}{mutualisme}. In een fosfaatrijke bodem
zou de plant het fosfaat zelf opnemen en kost de schimmel \qty{15}{\%} voor
niets: \omterm{def:b1:species-interactions:types}{parasitisme}. Proef: kweek planten met en zonder de schimmel bij een
reeks fosfaatniveaus en meet de biomassa; de verbintenis is mutualistisch
waar de planten met \omterm{def:b1:plant-water-minerals:mycorrhiza}{mycorrhiza} groter worden, parasitair waar zij kleiner
blijven.
\end{solution}

\begin{solution}{exo:b1:species-interactions:12}
Uitrasteringen: omheinde percelen langs de rivieren waar wapiti's niet in
kunnen, gekoppeld aan open percelen, herhaald in verscheidene dalen met en
zonder wolvenroedels, en jarenlang vóór en na de herintroductie gemeten;
noteer de hoogte en de bedekking van de wilgen, de dichtheid van de wapiti's
en de tijd die zij in elk perceel doorbrengen (sporen, camera's), de jacht en
de sneeuwdiepte. De cascade voorspelt dat de wilgen zich in de open percelen
alleen daar herstellen waar wolven zijn, en even snel als binnen de
uitrasteringen; herstellen de wilgen zich met en zonder wolven even goed, of
volgen zij de sneeuw en de jacht, dan is de cascade weerlegd of gedeeld met
die factoren.
\end{solution}

\begin{solution}{pb:b1:species-interactions:1}
\textbf{1.} A: $N_A + 1.4 N_B = 200$, snijpunten 200 ($N_A$) en 143
($N_B$). B: $N_B + 0.9 N_A = 130$, snijpunten 130 ($N_B$) en 144
($N_A$).
\textbf{2.} De snijpunten (200 tegen 144 op de $N_A$-as, 143 tegen 130 op de
$N_B$-as) leggen de \omterm{prop:b1:species-interactions:lotka}{isocline} van A geheel buiten die van B: het oplossen van
$N_A = 200 - 1.4 N_B$ met $N_B = 130 - 0.9 N_A$ geeft $N_A = -69$, zodat de
lijnen elkaar buiten het positieve kwadrant ontmoeten. A wint vanaf elk
beginpunt: overal waar B nog kan groeien, groeit A ook en houdt zij het
langer vol.
\textbf{3.} $T_A = 0.87$ dagen, $T_B = 1.26$ dagen: A, de snellere
groeier.
\textbf{4.} A: $N_A + 0.3 N_C = 200$ (200, 667); C: $N_C + 0.3 N_A =
100$ (100, 333). Elke soort begrenst zichzelf sterker dan de andere
($0.3 < 200/100$ en $0.3 < 100/200 = 0.5$): stabiel samenleven.
\textbf{5.} Een bos is niet constant (seizoenen en verstoringen zetten de
wedloop terug voordat zij beslist is) en heeft vele hulpbronnen en plaatsen,
zodat soorten nissen verdelen in plaats van er één te delen.
\textbf{6.} $N_A = 200 - 0.3 N_C$, $N_C = 100 - 0.3 N_A = 100 - 60 +
0.09 N_C$, dus $N_C = 44$ en $N_A = 187$.
\textbf{7.} Stijgend en afvlakkend naar ongeveer 10: type II.
\textbf{8.} $(1/N, 1/f)$: $(0.2, 0.5)$, $(0.1, 0.303)$, $(0.05, 0.2)$,
$(0.025, 0.149)$, $(0.0125, 0.125)$: een rechte lijn.
\textbf{9.} Helling $(0.5 - 0.125)/(0.2 - 0.0125) = 2.0 = 1/a$, dus $a =
\qty{0.5}{m^2/d}$; snijpunt $0.5 - 2\times 0.2 = 0.1$ dag $= h$.
\textbf{10.} Maximaal $1/h = 10$ prooien per dag; halve verzadiging bij $N =
1/(ah) = 20$ per vierkante meter.
\textbf{11.} Aandeel hanteren $= f h = f/10$: meer dan de helft wanneer $f
> 5$, dat wil zeggen bij 40 en 80 per vierkante meter (5,0 bij 20 is precies de helft).
\textbf{12.} De groei van de prooi $0.1 N (1 - N/200)$ is gelijk aan het
verbruik $0.5N/(1 + 0.05N)$: $0.1(1 - N/200) = 0.5/(1 + 0.05 N)$, dat wil zeggen
$(1 - N/200)(1 + 0.05N) = 5$: $1 + 0.05N - N/200 - N^2/4000 = 5$, dus
$N^2 - 180 N + 16\,000 = 0$, $N = 90 \pm \sqrt{8100 - 16\,000}$: geen
reële \omterm{def:b1:flowering-plant-organization:organs}{wortel} --- één kever per vierkante meter eet meer dan de prooi ooit kan
voortbrengen (de grootste groei is $rK/4 = 5$ per dag, bij $N = 100$, waar
de kever er 8,3 eet) en drijft haar tot uitsterven. Er is geen evenwicht; de
prooi houdt alleen stand als er minder kevers zijn of als zij toevluchten heeft.
\textbf{13.} De \omterm{def:b1:species-interactions:functional}{hanteertijd} halveert voor de tweede prooi ($h' = 0.05$),
zodat de kever er bij verzadiging tot 20 per dag van eet: de totale opname
bij een hoge dichtheid verdubbelt ruwweg als hij op de snellere prooi
overschakelt.
\textbf{14.} $H_{15} = \ln 15 = 2.71$. Acht soorten: $-(0.85\ln 0.85
+ 7\times 0.0214\ln 0.0214) = 0.138 + 0.576 = 0.71$.
\textbf{15.} Met een factor $2.71/0.71 = 3.8$.
\textbf{16.} De $-$ van de zeester op de mossel hield de $-$ van de mossel
op haar mededingers zwak; neem de eerste weg en de mossel bereikt de
dichtheid waarbij haar \omterm{def:b1:species-interactions:types}{mededinging} de rest uitsluit.
\textbf{17.} Biomassa meet hoeveel van de energiestroom een soort
vasthoudt, niet wat haar wisselwerkingen doen; de uitwerking van een
\omterm{def:b1:species-interactions:keystone}{sluitsteensoort} loopt via de soorten die zij in toom houdt, en die hebben
wel een grote biomassa.
\textbf{18.} Het wegnemen van de zeester zou een zeldzame soort hebben
vrijgelaten die de rots niet had kunnen overnemen: de door mosselen
overheerste gemeenschap zou weinig zijn veranderd, en de zeester zou geen
\omterm{def:b1:species-interactions:keystone}{sluitsteensoort} zijn geweest.
\textbf{19.} Kooien die de zeester van kleine plekken weren, met open kooien
als controle voor het effect van de kooi zelf: nemen de mosselen alleen in de
gesloten kooien de macht over, dan is de \omterm{def:b1:species-interactions:types}{predatie} op mosselen het mechanisme;
kooien met de zeester erin maar met de mosselen met de hand verwijderd zouden
uitwijzen of de andere prooien van de zeester ertoe doen.
\textbf{20.} Winst $\qty{40}{mg}$, kosten $8\times 0.5 = \qty{4}{mg}$: netto
\qty{36}{mg} $=$ \qty{0.58}{kJ}.
\textbf{21.} \qty{40}{\micro g} van \qty{4}{mg}: \qty{1}{\%} van het
dagelijkse fotosyntheseproduct van de bloem koopt ongeveer drie
stuifmeeloverdrachten, de voortplanting van de plant.
\textbf{22.} Elk neemt iets van de ander, maar elk wint meer dan het
verliest: het teken is dat van de netto uitwerking, niet van de uitbuiting.
\textbf{23.} Rover $+$, plant $-$ (nectar genomen, geen bestuiving): een
\omterm{def:b1:species-interactions:types}{parasitisme} van het \omterm{def:b1:species-interactions:types}{mutualisme}; de bij vindt lege bloemen en wint minder,
een indirecte $-$.
\textbf{24.} Plant $\leftrightarrow$ bij $+/+$; rover $\to$ plant $-$,
rover $\to$ bij $-$; vogel $\to$ rover $-$. Twee minnen: de vogel heeft een
indirecte $+$ op de plant en op de bij.
\textbf{25.} $a = \qty{0.5}{m^2/d}$, $h = \qty{0.1}{d}$ (\qty{2.4}{h}),
maximaal \num{10} prooien per dag.
\end{solution}
